% Solutions to the exercises of Chapter 30.
\begin{solutions}
  \item The core of a minimal guard:
\begin{lstlisting}
inline int g_allocations = 0;                  // reset before, check after
void* operator new(std::size_t n) {
  ++g_allocations;
  if (void* p = std::malloc(n)) return p;
  throw std::bad_alloc();
}
void operator delete(void* p) noexcept { std::free(p); }
void operator delete(void* p, std::size_t) noexcept { std::free(p); }
\end{lstlisting}
        Reset the counter immediately before the code under test and read it immediately after, because the test framework allocates too. Measured
        with clang/libc++:
        \begin{itemize}
          \item four \code{push\_back} calls into a vector with \code{reserve(4)}: 0 allocations; the fifth: 1 (the reallocation);
          \item resizing a fixed-max $75\times75$ matrix from $15\times15$ to $75\times75$: 0, because the storage is inline;
          \item a dynamic \code{Eigen::MatrixXd(10, 10)}: also 0. This is not because it does not allocate, but because Eigen allocates with \code{std::malloc}
                through its own aligned allocator, which never calls \code{operator new}.
        \end{itemize}
        With \code{EIGEN\_RUNTIME\_NO\_MALLOC} and \code{set\_is\_malloc\_allowed(false)}, the same \code{MatrixXd} aborts with Eigen's assertion ``heap allocation is forbidden''. The guard and
        the flag are therefore both needed. Since the flag is defined only in the \code{debug-sanitize} preset, Eigen allocations are caught only there. Defining it for the
        test targets of both presets would close that gap; it is a design decision to record in \code{PROGRESS.md}. A complete guard also overrides the aligned
        forms \code{operator new(std::size\_t, std::align\_val\_t)}, which the library implements without calling the plain one. The C++ experiments of these
        solutions are in \code{scripts/solutions/cpp}, with \code{build\_and\_run.ps1} to build and run them in both configurations.

  \item Five candidates, each assigned afterwards to a concrete type and built in both presets:
        \begin{itemize}
          \item \code{auto e = R * (a + b);} gives correct numbers in both builds. The sum is a nested expression, which Eigen stores by value, so nothing dangles.
          \item \code{auto e = (a + b).eval().transpose();} prints the right numbers in release, by luck. ASan reports \code{stack-use-after-scope}.
          \item \code{auto e = make().head<2>();}, with \code{make()} returning a \code{Vector3d} by value, behaves the same way: right by luck in release, reported by ASan.
          \item \code{auto e = (X * 2.0).eval().block(0, 0, 2, 2);} with a dynamic \code{X} prints garbage ($10^{-311}$ entries) in release without crashing.
                ASan reports \code{stack-use-after-scope}.
          \item \code{auto e = A * v;} followed by a change to \code{A} produces the product with the \emph{new} \code{A} in both builds, and no sanitizer reports anything.
        \end{itemize}
        Dangling references are undefined behavior: the optimized build may print the right numbers because
        the dead temporary's stack slot has not been reused yet, or garbage, as in the dynamic case. AddressSanitizer reports them reliably, as a read of a variable
        whose scope has ended. The stale expression is not undefined behavior at all, since every access is to valid memory. The expression is simply evaluated
        later than the reader assumed, so no sanitizer can report it. Only the rule ``no \code{auto} with Eigen'' prevents it.

  \item Ten thousand propagation steps of a $75\times75$ covariance (the $15\times15$ block with $\Phi P\Phi\T + Q$, the $15\times60$ cross block with $\Phi P_{IC}$), per step:
        \begin{center}\small
        \begin{tabular}{@{}l c c@{}}
        \toprule
        implementation & release & debug-sanitize \\
        \midrule
        fixed-max, member workspaces & \SI{3.0}{\micro\second} & \SI{500}{\micro\second} \\
        dynamic, temporaries in expressions & \SI{3.3}{\micro\second} & \SI{525}{\micro\second} \\
        dynamic, preallocated workspaces & \SI{3.3}{\micro\second} & \SI{505}{\micro\second} \\
        \bottomrule
        \end{tabular}
        \end{center}
        In release, the arithmetic (about $4\times10^4$ flops) dominates. The three small heap allocations per step of the naive dynamic version cost a few hundred
        nanoseconds, and fixed sizes let Eigen unroll the $15\times15$ parts. That gives a 10\% difference, not an order of magnitude. The dynamic version with workspaces does not allocate
        either (it ran with Eigen allocations forbidden). The case for fixed-max matrices is therefore not average speed but the \emph{guarantee}: the storage is
        inline, its bound is checked, and no forgotten \code{.noalias()} or temporary can allocate in the hot path. That is what makes the worst case predictable. The
        sanitizer build is about 170 times slower. At \code{-O0} none of Eigen's small template functions is inlined, and ASan checks every load and store. Use it to
        prove correctness (both allocation-free versions ran with \code{malloc} forbidden), never to measure speed.

  \item With the right Jacobian $\Jr$ of \cref{ch:lie}, tested against central differences of $\Log(\Exp(\phi)\T\Exp(\phi + \delta))$ ($h = 10^{-6}$, so the correct version agrees to $10^{-10}$), the
        largest error of three broken versions at four test points is:
        \begin{center}\small
        \begin{tabular}{@{}l c c c c@{}}
        \toprule
        breakage & $\phi = 0$ & $\|\phi\| = 0.05$ & $\|\phi\| = 0.5$ & $\|\phi\| = 2.5$ \\
        \midrule
        sign of the $\sx{\phi}$ term flipped (this is $\Jl$) & 0 & $4\times10^{-2}$ & $0.4$ & $1.1$ \\
        second-order term dropped & 0 & $4\times10^{-4}$ & $4\times10^{-2}$ & $0.65$ \\
        Taylor branch used below $0.1$~rad & 0 & $4\times10^{-4}$ & $8\times10^{-11}$ & $9\times10^{-11}$ \\
        \bottomrule
        \end{tabular}
        \end{center}
        A test at $\phi = 0$ detects nothing, since every version equals $I$ there. A test at a single generic angle misses the wrong branch threshold. A loose tolerance
        ($10^{-3}$, say) misses the dropped term at small angles. Hence the rules: several random points spanning the regimes (tiny, inside every branch, generic, near
        $\pi$), and a tolerance derived from the finite-difference error rather than chosen for comfort.

  \item Reading \code{a[4]} of a \code{std::array<int, 4>} and computing \code{INT\_MAX + 1}:
        \begin{itemize}
          \item \emph{release}: the read returned 1634890337 (whatever lay next on the stack) and the addition printed $-2147483648$. No crash and no message.
                Neither result is guaranteed: both are undefined behavior, and the optimizer may, for example, fold \code{x + 1 > x} to \code{true};
          \item \emph{debug-sanitize}: ASan stops at the read with \code{stack-buffer-overflow ... READ of size 4} and the file and line of the access. UBSan stops at the
                addition with \code{runtime error: signed integer overflow: 2147483647 + 1 cannot be represented in type 'int'}. With \code{-fno-sanitize-recover=all}
                both abort, so a test that triggers them fails.
        \end{itemize}
        Note that \code{std::array::operator[]} does not check bounds. ASan caught the read because the array lives on the stack between redzones. Compiling with libc++'s
        hardening (\code{-D\_LIBCPP\_HARDENING\_MODE=\_LIBCPP\_HARDENING\_MODE\_FAST}) adds a cheap bounds check that traps even in an optimized build: the process stopped with an
        illegal-instruction exception.

  \item Binding a $3\times3$ block of a row-major matrix to \code{Eigen::Ref<const Eigen::MatrixXd>} compiles and gives the right result, but \code{Ref} cannot describe a row-major block
        with column-major strides, so it copies the block into an internal \code{MatrixXd}. The operator-new guard counts 0 allocations, because the copy is an Eigen
        allocation through \code{malloc}. With \code{set\_is\_malloc\_allowed(false)} the call aborts with Eigen's assertion. A column-major block binds without any copy (it passes with
        allocations forbidden). With a writable \code{Eigen::Ref<Eigen::MatrixXd>}, the call does not compile: \code{no matching function for call to 'scale'}. A writable reference cannot
        silently redirect writes into a copy. Keep one storage order throughout the core (column-major, Eigen's default), or write the \code{Ref} type for the layout
        actually passed.

  \item With \code{PoseEvent} added to the variant and no lambda for it, libc++ stops with
\begin{lstlisting}[language={}, basicstyle=\ttfamily\footnotesize]
error: static assertion failed due to requirement
  'is_invocable_v<Overloaded<...>, const PoseEvent &>':
  std::visit requires the visitor to be exhaustive.
\end{lstlisting}
        This is followed by notes leading back to the \code{std::visit} call. Every dispatch site that forgot the new event type fails to compile, which is the point of the
        design. The one way to defeat it is a catch-all generic lambda, \code{[](const auto\&) \{\}}, which silently swallows new alternatives. Do not put one in the
        event dispatch.
\end{solutions}
